% !TEX root = ../main.tex
\chapter{一元函数微分学}
这一章节，我们将进入微分学，从高中的导数出发，我们假定你有良好的导数基础，因此我们会跳过大量前置知识，并默认你已了解。


\section{微分}

当自变量 $x$ 发生一个微小变化 $\Delta x$ 时，函数值 $y = f(x)$ 的变化量 $\Delta y$ 能否用 $\Delta x$ 的线性表达式来近似？也就是说，是否可以构造一个形似$y=kx$的表达式来近似$\Delta y$？更关键的是，这个近似的误差是否比 $\Delta x$ 本身小得多？由此我们引入可微的概念：

\subsection{可微与微分}

\begin{definition}{可微的定义}{}
  设函数 $y = f(x)$ 在点 $x_0$ 的某个邻域内有定义。记自变量的增量为 $\Delta x$，函数值的增量为$\Delta y$，则：
  \[
  \Delta y = f(x_0 + \Delta x) - f(x_0)
  \]
  若存在一个与 $\Delta x$ 无关的常数 $A$，使得：
\begin{center}
    当$\Delta x \to 0$时，$\Delta y = A \cdot \Delta x + o(\Delta x)$
\end{center}
  我们就说 $f(x)$ 在 $x_0$ 处\textbf{可微}。
\end{definition}

这个式子的意思是，函数值的变化量 $\Delta y$ 可以分解为两部分。第一部分 $A \cdot \Delta x$ 是 $\Delta x$ 的线性函数，称为\textbf{线性主部}。第二部分 $o(\Delta x)$ 是比 $\Delta x$ 更高阶的无穷小，它除以 $\Delta x$ 的商趋于 $0$。也就是说，当 $\Delta x$ 足够小时，$\Delta y$ 的主体就是 $A \Delta x$，$o(\Delta x)$比$\Delta x$ 趋于$0$的速度更快，误差可以忽略不计。

由此，我们可以引出微分的定义：

\begin{definition}{微分}{}
若 $y=f(x)$ 在 $x_0$ 处可微，则称线性主部 $A \Delta x$ 为 $y$ 在点 $x_0$ 处的\textbf{微分}，记作
\[
dy = A \Delta x
\]
对于自变量 $x$ 本身，$dx = \Delta x$。因此$y$ 对 $x$ 的微分也常记作
\[
dy = A dx
\]
\end{definition}

这里的常数 $A$ 是什么？我们把 $\Delta y = A \cdot \Delta x + o(\Delta x)$ 两边同时除以 $\Delta x$：
\[
\frac{\Delta y}{\Delta x} = A + \frac{o(\Delta x)}{\Delta x}
\]

又因为 $\Delta x \to 0$，所以$\frac{o(\Delta x)}{\Delta x}$趋于$0$，于是：
\[
\lim_{\Delta x \to 0} \frac{\Delta y}{\Delta x} = A
\]

也就是说，$A$ 不是别的，正是我们熟悉的导数。

所以可微等价于可导，且 $A = f'(x_0)$。微分和导数的关系可以写成：

\begin{theorem}{微分和导数的关系}{}
  若 $y=f(x)$ 在 $x_0$ 处可微，则：
  \[
  dy = f'(x_0)\, dx
  \]

  反之亦成立。
\end{theorem}

这个关系实则暗含了两层意思：可微必可导，可导必可微。

因此，莱布尼茨用记号 $\frac{dy}{dx}$ 表示 $f'(x)$，这表示导数是两个微小量 $dy$ 和 $dx$ 的比值，只不过这个比值在取极限后从近似变成了精确。所以，导数也称作“微商”（不是微信电商的意思）。

从几何上看，$dy = f'(x_0)\,\Delta x$ 是切线上的增量，也就是沿切线方向走 $\Delta x$ 所上升的高度，而 $\Delta y$ 是曲线上的真实增量。当 $\Delta x \to 0$ 时，两者的差 $|\Delta y - dy|$ 是比 $\Delta x$ 更高阶的无穷小。这也体现了微分“以直代曲”的数学思想，在极小范围内，函数的曲线和它的切线几乎重合。

\begin{myexample}
  设 $y = x^3$，讨论函数在 $x = 1$ 处的可微性，并求 $\Delta x = 0.01$ 时的微分和真实增量。

  先求 $A$。计算极限：
  \[
  \lim_{\Delta x \to 0} \frac{(1+\Delta x)^3 - 1^3}{\Delta x}
  = \lim_{\Delta x \to 0} \frac{3\Delta x + 3(\Delta x)^2 + (\Delta x)^3}{\Delta x}
  = 3
  \]

  因此 $f(x) = x^3$ 在 $x=1$ 处可微，且 $A = 3$。

  微分：$dy = A \Delta x = 3 \cdot 0.01 = 0.03$。

  真实增量：$\Delta y = (1.01)^3 - 1^3 = 1.030301 - 1 = 0.030301$。

  误差： $|\Delta y - dy| = 0.000301$，是 $(\Delta x)^2$ 量级的，确实比 $\Delta x$ 小得多。随着 $\Delta x$ 减小，误差会以更快的速度趋近于$0$。
\end{myexample}

\subsection{微分运算法则}

我们已经知道微分与导数的关系 $dy = f'(x)\,dx$，因此求微分的本质仍然是求导。但微分有自己独立的运算法则，它们在形式上比导数法则更对称，在积分学中尤其方便。我们默认你已经掌握了基本初等函数的求导公式，此处将它们转写成微分形式：

\begin{theorem}{微分公式}{}
  \[
  \begin{aligned}
  &\text{(1)}\ d(C) = 0
  &&\text{(2)}\ d(x^\mu) = \mu x^{\mu-1}\, dx \\[4pt]
  &\text{(3)}\ d(\sin x) = \cos x\, dx
  &&\text{(4)}\ d(\cos x) = -\sin x\, dx \\[4pt]
  &\text{(5)}\ d(\tan x) = \sec^2 x\, dx
  &&\text{(6)}\ d(\cot x) = -\csc^2 x\, dx \\[4pt]
  &\text{(7)}\ d(\sec x) = \sec x \tan x\, dx
  &&\text{(8)}\ d(\csc x) = -\csc x \cot x\, dx \\[4pt]
  &\text{(9)}\ d(e^x) = e^x\, dx
  &&\text{(10)}\ d(a^x) = a^x \ln a\, dx \\[4pt]
  &\text{(11)}\ d(\ln x) = \frac{1}{x}\, dx
  &&\text{(12)}\ d(\log_a x) = \frac{1}{x \ln a}\, dx \\[4pt]
  &\text{(13)}\ d(\arcsin x) = \frac{1}{\sqrt{1-x^2}}\, dx
  &&\text{(14)}\ d(\arccos x) = -\frac{1}{\sqrt{1-x^2}}\, dx \\[4pt]
  &\text{(15)}\ d(\arctan x) = \frac{1}{1+x^2}\, dx
  \end{aligned}
  \]
\end{theorem}

这些公式不需要死记，只需要将任何一个求导公式两边乘以 $dx$，就得到了对应的微分公式。

接下来是微分的四则运算法则。设 $u = u(x)$，$v = v(x)$ 都可微，则有：

\begin{theorem}{微分的四则运算法则}{}
\[
\begin{aligned}
&\text{(1)}\ d(u \pm v) = du \pm dv \\[4pt]
&\text{(2)}\ d(cu) = c\, du \quad (c\text{ 为常数}) \\[4pt]
&\text{(3)}\ d(uv) = u\, dv + v\, du \\[4pt]
&\text{(4)}\ d(\frac{u}{v}) = \frac{v\, du - u\, dv}{v^2} \quad (v \neq 0)
\end{aligned}
\]
\end{theorem}

这些法则与导数法则一一对应，只是换了一个写法。

也就是说，我们怎么求导，就怎么求微分，最后别忘记加上$dx$和$dy$就行。

\begin{myexample}
  设 $y = x^2 \sin x$，求 $dy$。

  用乘法法则：令 $u = x^2$，$v = \sin x$。则
  \(du = 2x\, dx\)，\( dv = \cos x\, dx\)。

  所以：
  \[
  dy = u\, dv + v\, du = x^2 \cdot \cos x\, dx + \sin x \cdot 2x\, dx = (x^2 \cos x + 2x \sin x)\, dx
  \]
\end{myexample}

\begin{myexample}
  设 $y = \dfrac{x}{1+x^2}$，求 $dy$。

  用除法法则：令 $u = x$，$v = 1+x^2$。则 $du = dx$，$dv = 2x\, dx$。

  所以：
  \[
  dy = \frac{v\, du - u\, dv}{v^2}
  = \frac{(1+x^2) \cdot dx - x \cdot 2x\, dx}{(1+x^2)^2}
  = \frac{1 + x^2 - 2x^2}{(1+x^2)^2}\, dx
  = \frac{1 - x^2}{(1+x^2)^2}\, dx
  \]
\end{myexample}

\section{导数}

导数在高中时期可谓是令人闻风丧胆。不过在这里，我们不会讨论导数难题，我们也不需要。和高中不同，高中导数学得深，而微分学的导数学得广。为什么广？往下看就知道了。

\subsection{可导与连续的关系}
我们假设你还记得导数的定义。不过不记得也没关系，现在帮你回忆一下：

\begin{definition}{导数}{}
设函数 $y = f(x)$ 在点 $x_0$ 的某个邻域内有定义。若极限
\[
\lim_{\Delta x \to 0} \frac{f(x_0 + \Delta x) - f(x_0)}{\Delta x}
\]
存在，则称此极限值为 $f(x)$ 在点 $x_0$ 处的\textbf{导数}，记作
\[
f'(x_0),\quad y'\big|_{x=x_0},\quad \frac{dy}{dx}\Big|_{x=x_0},\quad \frac{df}{dx}\Big|_{x=x_0}
\]
此时称 $f(x)$ 在点 $x_0$ 处\textbf{可导}。
\end{definition}

导数的几何意义是曲线在 $(x_0, f(x_0))$ 处的切线斜率，物理意义是瞬时变化率。

关于导数和连续两个概念，我们有这样的关系：

\begin{theorem}{可导必连续}{}
若 $f(x)$ 在点 $x_0$ 处可导，则 $f(x)$ 在点 $x_0$ 处必连续。
\end{theorem}

要证明这个结论并不难，读者可尝试自行证明。这里给出证明的大致思路：

我们知道，若 $f(x)$ 在点 $x_0$ 处可导，则 $f'(x_0)$ 存在，那么有：
\[
\lim_{\Delta x \to 0} \frac{f(x_0 + \Delta x) - f(x_0)}{\Delta x} = f'(x_0)
\]

因为 $\Delta y = f(x_0 + \Delta x) - f(x_0)$，那么：
\[
\lim_{\Delta x \to 0} \Delta y
= \lim_{\Delta x \to 0} ( \frac{\Delta y}{\Delta x} \cdot \Delta x )
= f'(x_0) \cdot 0 = 0
\]

而这正是连续性的增量语言 $\lim_{\Delta x \to 0} \Delta y = 0$，所以 $f(x)$ 在 $x_0$ 处连续。

直观上理解，导数存在意味着 $\frac{\Delta y}{\Delta x}$ 趋近于一个有限值，而当 $\Delta x$趋近于$0$时，$\Delta y$ 也必须趋近于 $0$，否则此时的导数就是$\frac{\text{常数}}{0}=\infty$，这不是一个有限值。

如果函数在某点不连续，那么它可能连极限都不存在，导数就更无从谈起了。所以\textbf{不连续一定不可导}。

我们能证明，\textbf{初等函数在其定义区间内都是连续的}。

$f(x)$ 在点 $x_0$ 处可导，则 $f(x)$ 在点 $x_0$ 处必连续。反过来成立吗？不成立，也就是说，函数连续却可能没有导数。

与极限和连续类似，我们先定义左导数和右导数：

\begin{definition}{左导数与右导数}{}
设函数 $y = f(x)$ 在点 $x_0$ 的某个左邻域或右邻域内有定义。

\textbf{左导数}：若极限
\[
\lim_{\Delta x \to 0^-} \frac{f(x_0 + \Delta x) - f(x_0)}{\Delta x}
\]
存在，则称此极限值为 $f(x)$ 在点 $x_0$ 处的\textbf{左导数}，记作 $f'_-(x_0)$。

\textbf{右导数}：若极限
\[
\lim_{\Delta x \to 0^+} \frac{f(x_0 + \Delta x) - f(x_0)}{\Delta x}
\]
存在，则称此极限值为 $f(x)$ 在点 $x_0$ 处的\textbf{右导数}，记作 $f'_+(x_0)$。

\end{definition}

左导数与右导数统称为单侧导数。与极限和连续类似，$f(x)$ 在 $x_0$ 处可导的充要条件是$f(x)$ 在 $x_0$ 处的左右导数都存在且相等。

\begin{theorem}{导数与单侧导数的关系}{}
  函数在一点可导的充要条件是函数在该点处的左右导数都存在且相等。即：
  \begin{center}
    若$f'_-(x_0)=f'_+(x_0)=f'(x_0)$，则$f(x)$在$x_0$处可导。
  \end{center}
  反之亦成立。
\end{theorem}

所以，即使函数连续，也无法推出左右导数相等，因此函数连续不一定可导。我们来看一个例子。

\begin{myexample}
  判断$f(x) = |x|$ 在 $x = 0$ 处的连续性和是否可导。

  连续性：由第一章已知 $\lim_{x \to 0} |x| = 0 = f(0)$，连续。

  可导性：先计算左右导数。
  
    左导数：$\quad \lim_{\Delta x \to 0^-} \frac{|0 + \Delta x| - 0}{\Delta x}
    = \lim_{\Delta x \to 0^-} \frac{-\Delta x}{\Delta x} = -1$

    右导数：$\quad \lim_{\Delta x \to 0^+} \frac{|0 + \Delta x| - 0}{\Delta x}
    = \lim_{\Delta x \to 0^+} \frac{\Delta x}{\Delta x} = 1$
  
  左右导数不相等，故导数不存在，因此$f(x) = |x|$ 在 $x = 0$ 处连续但不可导。

\end{myexample}

那么$f(x) = |x|$ 在 $x = 0$ 处可微吗？刚刚我们说到，可微必可导，可导必可微，所以$f(x) = |x|$ 在 $x = 0$ 处不可微。

几何上看，$f(x)=|x|$ 在 $x=0$ 处有一个“尖角”，左边的切线斜率和右边的切线斜率不一样，画不出一条唯一的切线。

总结一下：

\textbf{连续是可导的必要条件，但不是充分条件}。可导必连续，连续不一定可导。具体是否可导，看左右导数是否相等。

\subsection{高阶导数}

导数的导数称为二阶导数，二阶导数的导数称为三阶导数，以此类推。一般地，$n$ 阶导数就是求导 $n$ 次的结果。记号上，$y = f(x)$ 的高阶导数有两种主流写法：
\begin{definition}{高阶导数的记号}{}
  若函数$y = f(x)$ 有任意阶导数，则其高阶导数记为：

  拉格朗日记号：
  \[
  f'(x),\quad f''(x),\quad f'''(x),\quad f^{(4)}(x),\quad \dots,\quad f^{(n)}(x)
  \]
  
  莱布尼茨记号：
  \[
  \frac{dy}{dx},\quad \frac{d^2y}{dx^2},\quad \frac{d^3y}{dx^3},\quad \dots,\quad \frac{d^ny}{dx^n}
  \]
  
\end{definition}

莱布尼茨记号的写法值得解释一下。$\frac{d^2y}{dx^2}$ 的意思是“对 $y$ 取两次微分除以 $dx$ 的平方”：
\[
\frac{d^2y}{dx^2} = \frac{d}{dx} (\frac{dy}{dx})
\]

分子上的 $d^2y$ 表示微分运算执行了两次，分母上的 $dx^2$ 实际上是 $(dx)^2$ 的简写。这个记号的优点是操作感极强：求两次导就是分子上两个 $d$，分母上两个 $dx$。

看几个简单例子：

\begin{myexample}
  求 $y = x^4$ 的各阶导数。

  $y' = 4x^3$，$y'' = 12x^2$，$y''' = 24x$，$y^{(4)} = 24$，$y^{(5)} = y^{(6)} = \cdots = 0$。

\end{myexample}

由此，我们有：
\begin{theorem}{$n$ 次多项式的高阶导数}{}
  一般地，$n$ 次多项式的 $n$ 阶导数为常数，$n+1$ 阶及以上的导数全为零。
\end{theorem}


\begin{myexample}
  求 $y = \sin x$ 的高阶导数。

  $\sin x$ 的导数呈周期为 $4$ 的循环：
  \[
  (\sin x)' = \cos x,\quad (\sin x)'' = -\sin x,\quad (\sin x)''' = -\cos x,\quad (\sin x)^{(4)} = \sin x
  \]
\end{myexample}

由此，我们还能推出$\cos x$的规律：
\begin{theorem}{三角函数的高阶导数}{}
  一般地，$(\sin x)^{(n)} = \sin (x + \frac{n\pi}{2})$，$(\cos x)^{(n)} = \cos (x + \frac{n\pi}{2})$。
\end{theorem}

\subsection{复合函数求导}

前面讨论的都是直接对 $x$ 求导，但是有时候函数往往是一层套一层的，例如 $y = \sin x^2$，它不是单纯的 $\sin$，而是 $\sin$ 里面又包了一个 $x^2$。这种函数的求导法则称为链式法则。

\begin{theorem}{链式法则}{}
设 $y = f(u)$ 在 $u_0$ 处可导，$u = g(x)$ 在 $x_0$ 处可导，且 $u_0 = g(x_0)$，则复合函数 $y = f(g(x))$ 在 $x_0$ 处可导，且
\[
\frac{dy}{dx} = f'(u_0) \cdot g'(x_0) = f'(g(x_0)) \cdot g'(x_0)
\]
即：
\[
\frac{dy}{dx} = \frac{dy}{du} \cdot \frac{du}{dx}
\]
\end{theorem}

$\frac{dy}{dx} = \frac{dy}{du} \cdot \frac{du}{dx}$的形式上就像是分数约分，$du$ 上下抵消。不过这只是一个记忆技巧，并非严格的数学推导，但它极好地捕捉了链式法则的直觉。

链式法则可以推广到任意多层的复合。三层复合时：
\[
\frac{dy}{dx} = \frac{dy}{du} \cdot \frac{du}{dv} \cdot \frac{dv}{dx}
\]

形象地说，链式法则就像剥洋葱，从最外层一层层往内求导，每一步的导数依次相乘。

我们来看最开始的例子：
\begin{myexample}
  求 $y = \sin x^2$ 的导数。

  令 $u = x^2$，则$y = \sin u$，则有：
  \[
  \frac{dy}{du} = \cos u, \qquad \frac{du}{dx} = 2x
  \]

  由链式法则：
  \[
  \frac{dy}{dx} = \cos u \cdot 2x = 2x \cos x^2
  \]
\end{myexample}

\begin{myexample}
  求 $y = e^{3x}$ 的导数。

  令 $u = 3x$，则$y = e^u$，则有：
  \[
  \frac{dy}{du} = e^u, \qquad \frac{du}{dx} = 3
  \]

  由链式法则：
  \[
  \frac{dy}{dx} = e^u \cdot 3 = 3e^{3x}
  \]
\end{myexample}

熟练之后可以省略设 $u$ 的步骤，直接用“外导乘内导”的口诀：先对整体求导，里面不变，再乘以内部的导数。例如：

\begin{myexample}
  求 $y = \ln(\cos x^2)$ 的导数。

  这是一个三层复合：最外层是 $\ln$，中间层是 $\cos$，最内层是 $x^2$。按照剥洋葱的步骤，从最外层一层层往里求导，每一步的导数依次相乘：

  最外层 $y = \ln u$，对 $u$ 求导得 $\frac{1}{u}$，里面 $u = \cos x^2$ 不变。

  中间层 $u = \cos v$，对 $v$ 求导得 $-\sin v$，里面 $v = x^2$ 不变。

  最内层 $v = x^2$，对 $x$ 求导得 $2x$。

  三步连乘：
  \[
  \frac{dy}{dx} = \frac{1}{\cos x^2} \cdot (-\sin x^2) \cdot 2x = -2x \cdot \frac{\sin x^2}{\cos x^2} = -2x \tan x^2
  \]

\end{myexample}

常犯的一个错误是只求了外层忘了内层。例如把 $\sin x^2$ 的导数写成 $\cos x^2$ 就停了。记住一个简单的自检方法，如果内层不是单纯的 $x$，链式法则就一定要出手，每一步链都不能掉。

\subsection{参数方程求导}

有时候，$x$ 和 $y$ 并不以 $y = f(x)$ 的形式直接关联，而是各自通过第三个变量 $t$ 来表达：
\(\begin{cases}
x = \varphi (t) \\
y = \phi (t)
\end{cases}\)，那么问题来了：怎么求 $y$ 对 $x$ 的导数？

微分给出了答案：

\begin{theorem}{参数方程求导法则}{}
  若由参数方程\(\begin{cases}
x = \varphi (t) \\
y = \phi (t)
\end{cases}\)确定的$y=f(x)$在对应区间上可导，则其导数为：

\[
\frac{dy}{dx} = \frac{\phi' (t)}{\varphi' (t)}
\]

\end{theorem}

由微分定义，我们有：

\[
dx=\varphi' (t)dt, \qquad dy=\phi' (t)dt
\]

两者相除，我们就得到了上述公式。

也就是说，由参数方程确定的$y=f(x)$，求 $y$ 对 $x$ 的导数，等于 $y$ 和 $x$ 各自对参数 $t$ 的导数之商。

\begin{myexample}
  已知圆的参数方程 \(\begin{cases}
x = \cos t \\
y = \sin t
\end{cases}\)，$t \neq k\pi$，求 $\frac{dy}{dx}$。

  先求 $x$ 和 $y$ 对 $t$ 的导数：
  \[
  \frac{dx}{dt} = -\sin t, \qquad \frac{dy}{dt} = \cos t
  \]

  代入公式：
  \[
  \frac{dy}{dx} = \frac{\cos t}{-\sin t} = -\cot t
  \]

\end{myexample}

如果需要求参数方程的二阶导数，按定义推导，由于\(
\frac{d^2y}{dx^2} = \frac{d}{dx} (\frac{dy}{dx}) = \frac{d(\frac{dy}{dx})}{dx}\)，所以先把$\frac{dy}{dx}$看成新的$y=\psi (t)$，得到$d\psi =\psi' (t) dt$，再求$dx$。此时$dx$还是不变，即$dx=\varphi' (t)dt$，所以此时的二阶导数应该是$\frac{d^2y}{dx^2} = \frac{d\psi }{dx} = \frac{\psi' (t)}{\varphi' (t)}$。

全部写开，我们得到：

\begin{theorem}{参数方程的二阶导数}{}
    若由参数方程\(\begin{cases}
x = \varphi (t) \\
y = \phi (t)
\end{cases}\)确定的$y=f(x)$在对应区间上有二阶导数，则其二阶导数为：
\[
\frac{d^2y}{dx^2} = \frac{\phi''(t)\,\varphi'(t) - \phi'(t)\,\varphi''(t)}{[\varphi'(t)]^3}
\]
\end{theorem}

注意二阶导不能直接写成$\frac{\phi ''(t)}{\varphi ''(t)}$，这是一个常见的错误，必须按定义先求一阶导再对 $t$ 求一次。

\begin{myexample}
  已知圆的参数方程 \(\begin{cases}
x = \cos t \\
y = \sin t
\end{cases}\)，$t \neq k\pi$，求$\frac{d^2 y}{dx^2}$。

刚刚我们已经求了$dx$和一阶导：
  \[
  dx = -\sin t dt, \qquad \frac{dy}{dx} = -\cot t
  \]

  下面用两种方法求 $\frac{d^2y}{dx^2}$：

  方法一：二阶导按定义为：
  \[
\frac{d^2y}{dx^2} = \frac{d}{dx} (\frac{dy}{dx}) = \frac{d(\frac{dy}{dx})}{dx}
  \]
  
  已求得 $\frac{dy}{dx} = -\cot t$，令$\psi (t)= \frac{dy}{dx} = -\cot t$，对 $t$ 求导：
  \[
  \frac{d\psi }{dt} = \csc^2 t = \frac{1}{\sin^2 t}
  \]

  代入得：
  \[
  \frac{d^2y}{dx^2} = \frac{d\psi }{dx} = \frac{d\psi }{dt} \cdot \frac{dt}{dx} = \frac{1}{\sin^2 t} \cdot \frac{1}{-\sin t} = -\frac{1}{\sin^3 t} = -\csc^3 t
  \]

  方法二：参数方程的二阶导公式为：
\[
\frac{d^2y}{dx^2} = \frac{\phi''(t)\,\varphi'(t) - \phi'(t)\,\varphi''(t)}{[\varphi'(t)]^3}
\]

  先求二阶导：
  \[\varphi ''(t) = (-\sin t)' = -\cos t, \qquad \phi ''(t) = (\cos t)' = -\sin t\]
  
  代入公式：
  \[
  \frac{d^2y}{dx^2}
  = \frac{(-\sin t)(-\sin t) - (\cos t)(-\cos t)}{(-\sin t)^3}
  = \frac{\sin^2 t + \cos^2 t}{-\sin^3 t}
  = -\frac{1}{\sin^3 t}
  = -\csc^3 t
  \]
  
\end{myexample}

两种方法结果一致。本质上没有区别。

\subsection{隐函数求导}

之前我们求导的函数都是 $y = f(x)$ 的形式，$y$ 被“显式”地写成 $x$ 的表达式。但实际中常有这样的方程：
\[
F(x, y) = 0
\]

它并不直接写出 $y = f(x)$，而是将 $x$ 和 $y$ 的关系统一放在一个式子里。这种形式确定的函数叫隐函数：

\begin{definition}{隐函数}{}
  若函数关系$y = f(x)$或$x = g(y)$由方程
  \[
  F(x, y) = 0
  \]
  所确定，则称该函数为\textbf{隐函数}。
\end{definition}

隐函数求导的思路非常直接：把方程 $F(x,y)=0$ 两边同时对 $x$ 求导，其中 $y$ 视为 $x$ 的函数 $y(x)$，碰到 $y$ 的项就用链式法则。用莱布尼茨记号写出来尤为清晰。

为什么遇到$y$要用链式法则呢？因为，任何含有$y$的项，本质上都是$x$的复合函数。例如，$\sin y$其实是$\sin y(x)$，$\ln y$其实是$\ln y(x)$。按照复合函数求导法则，先对外层函数求导，里面不变，再乘以内部函数$y(x)$的导数，也就是$\frac{dy}{dx}$。

\begin{theorem}{隐函数求导法则}{}
  若函数关系$y = f(x)$由方程
  \[
  F(x, y) = 0
  \]
  所确定，则求 $y$ 对 $x$ 的导数时，方程两边同时对$x$求导，其中 $y$ 视为 $x$ 的函数 $y=y(x)$，最后解出$\frac{dy}{dx}$。
\end{theorem}

\begin{myexample}
  设 $x^2 + y^2 = 1$（$y \geq 0$），求 $\frac{dy}{dx}$。

  两边对 $x$ 求导：
  \[
  \frac{d}{dx}(x^2) + \frac{d}{dx}(y^2) = \frac{d}{dx}(1)
  \]

  左边第一项是 $2x$，右边是 $0$。左边第二项中 $y^2$ 是 $y$ 的函数，$y$ 又是 $x$ 的函数，由链式法则：
  \[
  \frac{d}{dx}(y^2) = 2y \cdot \frac{dy}{dx}
  \]
  
  因此：
  \[
  2x + 2y \frac{dy}{dx} = 0
  \]

  解得：
  \[
  \frac{dy}{dx} = -\frac{x}{y}
  \]

\end{myexample}

结果中同时包含 $x$ 和 $y$，这在隐函数求导中是很正常的，我们不一定要把 $y$ 用 $x$ 表达出来。

  我们可以验证，由 $x^2+y^2=1$ 解得 $y = \sqrt{1-x^2}$，直接求导得 $\frac{dy}{dx} = \frac{-x}{\sqrt{1-x^2}} = -\frac{x}{y}$，这与我们求出来的结果一致。

\begin{myexample}
  设方程 $e^y + xy = x$ 确定了隐函数 $y = y(x)$，求 $\frac{d^2 y}{dx^2}$。
  
  左边 $e^y$ 对 $x$ 求导：
  \[\frac{d}{dx}(e^y) = e^y \cdot \frac{dy}{dx}\]
  
  $xy$ 是乘法，求导要注意：
  \[\frac{d}{dx}(xy) = 1 \cdot y + x \cdot \frac{dy}{dx}\]
  
  右边 $x$ 求导得 $1$。于是：
  \[
  e^y \frac{dy}{dx} + y + x \frac{dy}{dx} = 1
  \]

  把含 $\frac{dy}{dx}$ 的项提到一边：
  \[
  (e^y + x) \frac{dy}{dx} = 1 - y
  \]

  解得：
  \[
  \frac{dy}{dx} = \frac{1 - y}{e^y + x}
  \]

  则二阶导数为：
  \[
  \frac{d^2 y}{dx^2} =\frac{d}{dx} (\frac{1 - y}{e^y + x})
  \]

  所以：
  \[
\begin{aligned}
\frac{d^2y}{dx^2} 
&= \frac{d}{dx}(\frac{1-y}{e^y+x}) \\[6pt]
&= \frac{ (1-y)'(e^y+x) - (1-y)(e^y+x)' }{ (e^y+x)^2 } \\[6pt]
&= \frac{ (-y')(e^y+x) - (1-y)(e^y\cdot y' + 1) }{ (e^y+x)^2 } \\[6pt]
&= \frac{ -y'(e^y+x) - (1-y)e^y y' - (1-y) }{ (e^y+x)^2 } \\[6pt]
&= \frac{ -y'[(e^y+x) + (1-y)e^y] - (1-y) }{ (e^y+x)^2 } \\[6pt]
&= \frac{ -y'(2e^y + x - ye^y) - (1-y) }{ (e^y+x)^2 }
\end{aligned}
\]

代入\(y' = \frac{dy}{dx} = \frac{1-y}{e^y+x}\)：
\[
\begin{aligned}
\frac{d^2y}{dx^2} 
&= \frac{ -\frac{1-y}{e^y+x}( 2e^y + x - ye^y ) - (1-y) }{ (e^y+x)^2 } \\[8pt]
&= \frac{ -(1-y)( 2e^y + x - ye^y ) - (1-y)(e^y+x) }{ (e^y+x)^3 } \\[6pt]
&= \frac{ -(1-y)[ 2e^y + x - ye^y + e^y + x ] }{ (e^y+x)^3 } \\[6pt]
&= -\frac{ (1-y)( 3e^y + 2x - ye^y ) }{ (e^y+x)^3 }
\end{aligned}
\]
\end{myexample}

隐函数求导的核心就一条：把 $y$ 当作 $x$ 的函数，每碰到 $y$ 的项就使用链式法则，乘以 $\frac{dy}{dx}$，然后从方程中把 $\frac{dy}{dx}$ 解出来。

实际上，方程\(F(x, y) = 0\)可看作关于$x,y$的二元函数，也就是说，在多元函数微分学里，我们有更强大的工具可以求隐函数的导数。


\subsection{一阶微分形式不变性}

学完各种各样的导数之后，我们可以推出微分的一个极其优雅的性质。设 $y = f(u)$，如果 $u$ 是自变量，微分自然是 $dy = f'(u)\,du$。但如果 $u$ 本身是另一个变量 $x$ 的函数 $u = g(x)$，那么复合函数 $y = f(g(x))$ 的微分会是什么形式？

用微分的定义来算。先写出 $y$ 对 $x$ 的微分，由链式法则，先对$f$求导，再乘以内部$g$的导数，也就是：
\[
dy = [f(g(x))]'\, dx = f'(g(x))\, g'(x)\, dx
\]

注意到 $g'(x)\, dx$ 正是 $u$ 作为 $x$ 的函数的微分，即 $du = g'(x)\, dx$。代入上式：
\[
dy = f'(u)\, du
\]

也就是说，无论 $u$ 是自变量还是中间变量，微分永远写成 $dy = f'(u)\, du$ 这个形式。

\begin{theorem}{一阶微分形式不变性}{}
设 $y = f(u)$ 可微。则不论 $u$ 是自变量还是中间变量，总有：
\[
dy = f'(u)\, du
\]
\end{theorem}

这个性质是微分的灵魂。在后面的换元积分法中你会反复看到这个性质的身影。它是连接微分学与积分学的第一座桥梁。

一阶微分形式不变性由莱布尼茨发现。不仅如此，莱布尼茨还发现了复合函数求导的链式法则，不过当时没有严谨证明，证明的工作最终由柯西等人完成。而隐函数求导法则的系统化，则在多元函数微分学建立之后，由柯西、雅可比等人发展为著名的隐函数定理，成为现代数学分析的基石。

关于导数，柯西在《分析教程》（1821）中将导数明确表述为：
\[
f'(x) = \lim_{\Delta x \to 0} \frac{f(x+\Delta x)-f(x)}{\Delta x}
\]

柯西还系统讨论了连续性与可导性的关系，证明了可导必连续。更进一步，德国数学家魏尔斯特拉斯构造了一个函数：
\[
f(x) = \sum_{n=0}^{\infty} a^n \cos(b^n \pi x)
\]

其中 \(0 < a < 1\)，\(b\) 是一个正奇数，且 \(ab > 1 + \frac{3\pi}{2}\)。

这个函数在每一处都连续，但在每一处都不可导。画出来大概是一团永远在抖动的锯齿，无论放多大，都看不到任何一段光滑的曲线。

魏尔斯特拉斯早年在波恩大学读法律与财政，成绩平平，随后转入数学，却在毕业考试中因病失败。此后长达十余年，他在偏远的普鲁士小城当中学教师，教授书法、地理和体育。所以，当你说别人的数学是体育老师教的时候，别忘了魏尔斯特拉斯曾经也是体育老师。